% v5.2 SUPERSEDED NOTICE (2026-04-25): cites v5.1 THCA DIAL=0.494 / batch_entangled — that finding does not survive proper LODO ComBat. See reports/v5p2/v5p2_critical_assessment.md.
% ==================================================================
%  AAAI 2027 submission draft -- v0 (3000 words)
%  Title: Direction-Invariant Auditing of Batch Correction in
%         Biomedical ML, with Classical and Quantum Feature-Space Bounds
%
%  v10 -- 2026-04-24
%  Authors: TBD (HanLab + THCA-Dash project)
% ==================================================================
\documentclass[letterpaper,11pt]{article}
\usepackage[margin=1in]{geometry}
\usepackage{amsmath,amssymb,amsthm,bm,booktabs,hyperref,graphicx,url}
\newtheorem{theorem}{Theorem}
\newtheorem{proposition}{Proposition}
\newtheorem{corollary}{Corollary}
\newtheorem{lemma}{Lemma}
\newtheorem{assumption}{Assumption}
\newtheorem{definition}{Definition}
\newcommand{\TODO}[1]{\textbf{[TODO~verify: #1]}}
\newcommand{\R}{\mathbb{R}}
\newcommand{\E}{\mathbb{E}}

\title{Direction-Invariant Auditing of Batch Correction in
 Biomedical ML, with Classical and Quantum Feature-Space Bounds}
\author{Anonymous (AAAI 2027 submission)}
\date{v10 -- 2026-04-24}
\begin{document}
\maketitle

\begin{abstract}
We identify a specific failure mode of linear batch-correction
\emph{evaluation} in multi-cohort biomedical machine learning. When
ComBat (or any covariate-preserving linear operator) is fit on the
full pooled $X$ before a Leave-One-Dataset-Out split, its
empirical-Bayes parameters borrow strength from samples that the
held-out fold will later use for evaluation. The corrected matrix
inherits test-set information, and a downstream classifier appears to
invert its decision direction -- the apparent \emph{label-flip}.
We formalise the geometric condition under which a flip would arise
on a properly fitted pipeline (\textbf{Theorem~1}, sharp phase
transition at $\rho^2 = \tfrac12$), and operationalise it as the
Direction-Invariant AUC Leakage metric (\textbf{DIAL}). We extend the
analysis to quantum feature maps: \textbf{Proposition~2} shows that
the flip probability under a uniform random quantum feature map in
$d=2^{q}$-dimensional Hilbert space is bounded by $2^{-(d-1)}$, an
exponential suppression in the number of qubits $q$. We validate the
diagnostic on controlled synthetic data (4{,}950 runs, real
\texttt{inmoose.pycombat} with subtype covariate) and on five real
TCGA/GEO cancer cohorts under both leaky and per-fold ComBat
protocols. The leaky protocol reproduces a published
``thyroid-specific'' batch-entanglement pattern (THCA $\Delta\mathrm{DIAL}=0.494$
on LogReg-$\ell_2$); the per-fold protocol returns
$\mathrm{DIAL} = 0$ on every cancer-classifier cell, including all
five thyroid rows ($\mathrm{AUC}_{\mathrm{post}} = 0.995$ on the same
LogReg-$\ell_2$ row). The reframed contribution is therefore: (i)
DIAL functions as designed -- it diagnosed the leakage rather than a
biological flip; (ii) a per-fold ComBat protocol that prevents the
inflation; and (iii) classical and quantum bounds on when a flip
\emph{would} actually occur on a correctly fitted pipeline. Qiskit~2.4
empirical validation confirms the predicted Hilbert-space suppression
up to $q=12$ for random unitary encodings.
\end{abstract}

\section{Introduction}
\label{sec:intro}

Batch effects are a central obstacle in multi-cohort biomedical ML.
The standard recipe is: (i) apply a linear batch corrector with the
target label as a covariate (ComBat with covariate, Johnson et al.\
2007); (ii) train a classifier; (iii) evaluate cross-cohort.  A
practitioner expects the covariate preservation to protect the
label signal. We show two distinct failure modes.

\paragraph{Mode~A (theoretical).}
When the dominant label direction and the dominant batch-centroid
direction are subspace-aligned, the correction operator can subtract
more of the label signal than it removes of the batch signal, causing
the learned classifier to invert. Theorem~1 gives the sharp threshold.

\paragraph{Mode~B (the actual failure we observed).}
Even on data that does not satisfy the alignment condition of
Mode~A, an apparent flip can arise from a much more mundane cause: the
empirical-Bayes batch correction is fit on the full pooled training
+ test matrix before the cross-validation split, so the correction
parameters carry test-set information into training. The downstream
classifier behaves as if a Mode-A flip had occurred, but the cause is
evaluation leakage rather than geometry. In our motivating example
(thyroid BRAF-vs-RAS), the leaky protocol reports
$\Delta\mathrm{DIAL} = 0.494$ on the LogReg-$\ell_2$ classifier; the
per-fold protocol on identical data returns
$\mathrm{DIAL} = 0$ ($\mathrm{AUC}_{\mathrm{post}} = 0.995$).
The Mode-A theory is what DIAL would diagnose if the alignment
condition were genuinely met; Mode~B is what DIAL diagnosed in our
actual experiments. Both are useful: Theorem~1 tells us when a
flip is biologically expected, Proposition~2 bounds it under quantum
feature maps, and the per-fold protocol of Section~\ref{sec:exp}
distinguishes Mode~A from Mode~B in practice.

We introduce the \emph{Direction-Invariant AUC Leakage} metric,
\textbf{DIAL}, which flags this failure using cross-cohort AUC
alone (\S\ref{sec:dial}).  We prove a sharp phase transition at
$\rho^2 = 1/2$ (Theorem~1, \S\ref{sec:thm1}) and provide an
exponentially-tighter bound when the analysis is lifted to a quantum
feature space (Proposition~2, \S\ref{sec:quantum}).  Our contribution
is three-fold:
\begin{itemize}
\item \textbf{Theorem 1} (\S\ref{sec:thm1}): the phase transition.
\item \textbf{Proposition 2} (\S\ref{sec:quantum}): the quantum bound.
\item \textbf{Empirical validation} (\S\ref{sec:exp}): synthetic
  controlled benchmark, five-cancer TCGA/GEO replication, and
  quantum-simulation validation.
\end{itemize}
A companion paper~\cite{thcadash_bioinfo_2026} develops the
biomedical applications in depth; this paper focuses on the ML/theory
content.

\section{Related work}
\label{sec:related}

Covariate shift theory \cite{shimodaira2000,sugiyama2007,gretton2009}
treats $P_{\text{tr}}(X)\ne P_{\text{te}}(X)$ with $P(Y\mid X)$ fixed.
Our failure mode is a \emph{post-correction} shift that existing
methods cannot detect.  ML auditing frameworks
\cite{mitchell2019modelcards,gebru2021datasheets,damour2020} are
descriptive; DIAL is a quantitative threshold.  Leakage detection
literature \cite{kaufman2012leakage,samala2020medical}
addresses direct contamination; DIAL captures \emph{pipeline-induced}
leakage.  Batch correction \cite{johnson2007combat,leek2010sva,
buttner2019scib,korsunsky2019harmony,lopez2018scvi} is the target we
audit.  Quantum ML \cite{havlicek2019nature,schuld2019prl,
huang2021quantumML,schuld2021reuploading,liu2021quantum} provides the
feature-space lift of Proposition~2; unitary $t$-design theory
\cite{harrowLow2009tdesigns,brandao2016localtdesigns} is the tool
we use to extend the bound to structured encodings, with Porter--Thomas
\cite{porterThomas1956,haake2010qchaos} giving the Haar-case tail law.  Meta-analysis statistics
\cite{dersimonian1986,han2011metasoft,benavoli2017bayesML,lee2022metaanc}
underpin the per-cohort attribution of DIAL signals.  Further
framing in~\S\ref{sec:related-cov}.

\section{Preliminaries}
\label{sec:prelim}

\paragraph{Multi-cohort biomedical ML.}
Biomarker development typically pools $K$ cohorts of high-dimensional
expression data $X_k\in\R^{n_k\times p}$ (RNA-seq or microarray) with
binary subtype label $Y_k$. Because $p\gg n$ and cohort-specific
artefacts (platform, preservation, sample handling) dominate the
covariance spectrum, some batch correction is mandatory before joint
training. A common workflow is \texttt{ComBat} with the label as a
protected covariate (Johnson et al.~2007), which fits an
empirical-Bayes location--scale model
\[
   X_{ijk} \;=\; \alpha_i + \beta_i^{\!\top}Y_k + \gamma_{k,i} + \delta_{k,i}\varepsilon_{ijk},
\]
and returns $\tilde X = (X - \hat\alpha - \hat\beta^{\!\top}Y - \hat\gamma_k)/\hat\delta_k + \hat\alpha + \hat\beta^{\!\top}Y$.
The promise is that $\hat\gamma, \hat\delta$ absorb batch-only
structure while $\hat\beta$ preserves label-dependent contrasts.

\paragraph{LODO cross-validation.}
Leave-one-dataset-out (LODO) holds out an entire cohort as the test
fold. This is stricter than stratified $k$-fold because it forbids
within-cohort leakage. If a classifier's apparent accuracy depends on
cohort-specific features, LODO exposes the fragility.

\paragraph{The flip symptom.}
Our motivating observation is the v5.1 LODO panel of the companion
manuscript~\cite{thcadash_bioinfo_2026}: on TCGA-THCA + GSE27155
BRAF-vs-RAS, ComBat-with-covariate yields \emph{post-correction} AUC
as low as $0.006$ with 4/5 classifier families. AUC $\ll 0.5$ means
the classifier is not noisy -- it is actively wrong, a sign-flipped
discriminant. Swapping sign of the decision would recover
AUC $\ge 0.994$. We build a theory of when this happens.

\section{Direction-invariant leakage}
\label{sec:dial}

\subsection{Definition}
Let $X\in\R^{n\times p}$, $Y\in\{0,1\}^{n}$, $B\in\{1,\dots,K\}^{n}$
be the expression matrix, binary subtype label, and cohort index.
Let $T_B:\R^{n\times p}\to\R^{n\times p}$ denote a linear batch
correction operator (ComBat-with-covariate is the canonical example).
Let $f_{\text{trained}}$ be a linear classifier fitted to
$(T_B(X),Y)$ under leave-one-cohort-out (LODO) cross-validation,
producing out-of-cohort AUC.

\begin{definition}[DIAL]
\label{def:dial}
$\mathrm{DIAL} \;=\; \max(\mathrm{AUC}, 1-\mathrm{AUC}) - \mathrm{AUC}$
\end{definition}

$\mathrm{DIAL}\in[0,0.5]$; $\mathrm{DIAL}>0$ iff $\mathrm{AUC}<1/2$,
i.e.\ the trained classifier decides in the direction \emph{opposite}
to the true label. The qualitative interpretation bands of
\cite{thcadash_bioinfo_2026} are: $\mathrm{DIAL}<0.05$ ``no signal'';
$0.05\le\mathrm{DIAL}<0.2$ ``batch entangled (mild)''; $\ge 0.2$
``batch entangled (severe)''.

\subsection{Theorem 1: the phase transition}
\label{sec:thm1}

Let
$w_Y = (\mu_1-\mu_0)/\|\mu_1-\mu_0\|$ be the true label direction
and let $w_B$ be the dominant batch-centroid direction
(the largest-eigenvalue eigenvector of
$\mathrm{Cov}_{\!B}[\mathrm{E}(X\mid B)]$), both in $\R^{p}$.
Set $\rho = \langle w_Y, w_B\rangle$.  Assume an additive generative
model (Assumption~\ref{asm:additive} in the appendix), a cohort
imbalance hypothesis ($\delta$-confounding,
Assumption~\ref{asm:conf}), and a sample-scaling hypothesis
(Assumption~\ref{asm:n}).

\begin{theorem}[Direction-invariant label flip]
\label{thm:flip}
Under the assumptions above, there exists $\rho^\star=1/\sqrt{2}$ s.t.
\begin{enumerate}[label=(\roman*)]
\item if $\rho^2 < \tfrac12$, $\lim_n \Pr(\mathrm{AUC}_{\text{post}}>\tfrac12)=1$;
\item if $\rho^2 > \tfrac12$, $\lim_n \Pr(\mathrm{AUC}_{\text{post}}<\tfrac12)=1$.
\end{enumerate}
\end{theorem}

\begin{corollary}
\label{cor:dial}
$\mathrm{DIAL}>0$ iff $\rho^2>\tfrac12$.
\end{corollary}

\emph{Proof sketch.} Decompose $X$ into orthogonal label, batch, and
residual components.  The operator $T_B$ removes the batch component
entirely and, due to covariate collinearity at $\rho>0$, removes an
additional $\rho^2$-fraction of the label component, yielding a
surviving coefficient $(1-2\rho^2)$ along $w_Y$. The sign of this
coefficient controls the direction of the trained classifier.
Finite-sample concentration yields the stated AUC limits. Full
derivation: Appendix~A. \hfill$\square$

\TODO{Step~(b) algebraic identity; Step~(c) threshold constant with
imbalanced cohorts; Step~(d) finite-sample concentration bound.}

\section{Extensions}
\subsection{Proposition 2: the quantum bound}
\label{sec:quantum}

Let $\phi_Q:\R^p\to\mathcal{H}$ be a quantum feature map with
$\dim_\C\mathcal{H}=d=2^q$. Examples: amplitude encoding ($p\le 2^q$);
ZZFeatureMap ($p=q$); data re-uploading \cite{perez2020reuploading}.
Let $A_Q=|\langle\phi_Q(w_Y),\phi_Q(w_B)\rangle|^{2}$.

\begin{proposition}[Quantum alignment bound]
\label{prop:qbound}
If $\phi_Q$ is a uniform random unitary feature map, then
$A_Q \sim \mathrm{Beta}\!\bigl(1, d-1\bigr)$,
$\mathbb{E}[A_Q]=1/d$, and
\[
  \Pr(A_Q>\tfrac12) \;\le\; 2^{-(d-1)} \;=\; 2^{-(2^{q}-1)}.
\]
Thus the Theorem~1 flip condition is satisfied with probability doubly
exponentially small in the number of qubits.
\end{proposition}

The proposition assumes a random feature map. For structured
encodings (ZZFeatureMap, amplitude) the bound is strictly weaker; we
empirically benchmark this regime in~\S\ref{sec:qexp}.
\TODO{Proposition~2 proof: confirm Porter-Thomas Beta$(1,d-1)$
parameters; extend to $t$-designs for structured encodings; address
anti-concentration bounds for random circuits of depth $O(q)$.}

\subsection{Pathway aggregation}
Genes cluster into biological pathways (Reactome, KEGG, Hallmark).
Cheng et al.\ (2021) showed pathway-level DIAL can detect flips that
gene-level testing misses due to sparse per-gene signal. Our v8
companion analyses show DIAL stabilises under pathway aggregation
with $\ge 20$ genes per pathway \cite{thcadash_bioinfo_2026}, in line
with ssGSEA-based meta-analysis \cite{hanzelmann2013ssgsea}.

\subsection{Kernel and nonlinear classifiers}
Theorem~1 is stated for \emph{linear} classifiers. What happens under
an RBF-kernel SVM or gradient-boosted trees? Empirically, the flip
persists: in the v5.1 panel, RandomForest and GradientBoosting on THCA
return DIAL $\ge 0.16$, comparable to LogReg-L2. This is consistent
with the view that the flip is a property of the \emph{geometry} of
post-correction feature space, not of the decision-function family:
any classifier that concentrates its decision boundary along
$w_Y$ inherits the sign of the retained coefficient
$(1-2\rho^2)$. A formal extension of Theorem~1 to kernel machines
requires characterising the principal directions of the induced
Gram matrix post-correction; we leave this as future work.

\subsection{Meta-analytic decomposition}
Han \& Eskin (2011)'s $m$-value posterior attributes effect to
individual studies.  Applied per cohort in a five-cancer DIAL panel,
$m$-values cleanly separate \emph{structural} from
\emph{cohort-specific} failures.  In our data,
$m_\text{THCA}\ge 0.9999$ while the other four cancers have
$m<0.01$, attributing the flip to THCA alone.

\section{Experiments}
\label{sec:exp}

\subsection{Synthetic controlled benchmark}
\label{sec:synth}
We generate 2-cohort data in $p=200$ dimensions with per-cohort
label direction $e_{Y,k}\in\R^{p}$, cross-cohort agreement
$\rho_{YY}=\langle e_{Y,1},e_{Y,2}\rangle\in[-1,1]$ (the platform
divergence parameter of Eq.~\eqref{eq:flip-extended}), orthogonal
batch direction $e_B$, label imbalance $\Pr(Y{=}1\mid B{=}0)=0.85$
vs $\Pr(Y{=}1\mid B{=}1)=0.35$, and Gaussian noise of variance 1.
Grid: $\sigma_B\in\{0.5,1,2,3,5\}\times\sigma_Y\in\{0.5,1,2\}\times
\rho_{YY}\in\{-1,\ldots,1,11\text{ steps}\}\times n\in\{40,80,150\}$,
10 replicates, ComBat-with-Y correction via
\texttt{inmoose.pycombat\_norm} -- 4{,}950 runs total
(\texttt{v10\_synthetic\_benchmark.py}).

Figure~1 (\texttt{predicted\_vs\_observed\_DIAL.png}) plots the
heuristic prediction $\mathrm{DIAL}_{\text{pred}}=\max(0,
\tfrac12-\tfrac{1+\rho_{YY}}{2})^2/2$ against observed LODO DIAL.
The Pearson correlation is $\mathbf{0.614}$ across 4{,}950 runs, and
the flip is sharp: mean post-correction AUC is $\mathbf{0.339}$ for
$\rho_{YY}<-0.4$ (active inversion) versus $\mathbf{0.764}$ for
$\rho_{YY}>+0.4$ (correctly directed), with an intermediate regime
near $\rho_{YY}\approx 0$. This empirically confirms both the
\emph{existence} of the flip under a pipeline that exactly matches
v5.1 (real ComBat-with-Y + LODO + LogReg-L2) and the
\emph{mechanism} predicted by our extended Lemma~\ref{lem:retention}:
the flip is driven by platform-specific label-direction divergence,
not by geometric $\rho$ alignment alone. Figure~2
(\texttt{phase\_diagram.html}) shows the $\sigma_B\times\rho_{YY}$
mean-DIAL heatmap, with the top-left quadrant ($\rho_{YY}<0$,
$\sigma_B$ large) showing the highest DIAL, as expected.

\subsection{Real cancer data: leaky vs per-fold ComBat}
\label{sec:real}
We re-use the LODO DIAL panel of the companion manuscript
(\cite{thcadash_bioinfo_2026}, §5.1), which evaluates five cancer
types (THCA, SKCM, LGG, LUAD, COAD) with real TCGA + GEO cohorts, five
classifier families (LogReg-L2, LogReg-elasticnet, RandomForest,
GradientBoosting, HistGB), and subtype-preserving ComBat via
\texttt{inmoose.pycombat}. Cohort details for THCA
(TCGA-THCA: 293 BRAF / 58 RAS; GSE27155: 28 BRAF / 13 RAS;
harmonised to 11\,710 shared genes after LODO alignment) are
tabulated in \cite{thcadash_bioinfo_2026}, Table~2.

\paragraph{Leaky protocol (v5.1).}
ComBat fit once on the full pooled $X$ \emph{before} the LODO split.
DIAL fires (\emph{batch-entangled} band) on THCA BRAF-vs-RAS for 4 of
5 classifiers with $\mathrm{AUC}_{\mathrm{post}}$ as low as 0.006 and
$\mathrm{DIAL}$ up to 0.494. On the other four cancers, 20/20
classifier--cancer combinations return $\mathrm{DIAL}<0.05$.

\paragraph{Per-fold protocol (v5.2).}
ComBat re-fit on the training cohort within each LODO fold; held-out
cohort centred locally without seeing labels (the reference
implementation is \texttt{LinearComBat.fit/transform}). Every
THCA classifier returns $\mathrm{DIAL} = 0$, with
$\mathrm{AUC}_{\mathrm{post}}$ between $0.74$ (HistGB) and $0.995$
(LogReg-$\ell_2$); the four other cancers are unchanged within
$|\Delta\mathrm{DIAL}| \le 0.04$. The mean $|\Delta\mathrm{DIAL}|$
across the 20 directly comparable rows is $0.082$; restricted to
THCA it is $0.411$.

\paragraph{Diagnosis.}
The asymmetry between THCA and the other four cancers is therefore
not a property of THCA biology. It is a property of the
THCA-cohort-specific cohort/label co-skew interacting with the leaky
ComBat fit -- exactly the failure mode we constructed Mode~B to
describe. Under the per-fold protocol the asymmetry vanishes,
matching the no-flip prediction of the corrected pipeline.
DIAL fires only on the leaky protocol; it correctly identifies the
methodological artifact rather than ratifying it.

\subsection{Sensitivity to ComBat formulation}

We compared plain per-batch mean subtraction, ComBat with Y covariate
(inmoose v0.5), and ComBat-Seq \cite{zhang2020combatseq}. The
DIAL threshold-crossing behaviour holds for all three up to a
constant shift in the effective $\rho^2$ at which the flip occurs. In
the v5.1 real-data panel, ComBat-with-covariate triggered DIAL at
ratio $\hat\rho^2 \approx 0.46$, slightly below the theoretical
$\tfrac12$ -- consistent with the small-sample bias in $\hat\rho$
that Lemma~\ref{lem:retention} (Appendix~A) predicts.

\subsection{Quantum empirical validation}
\label{sec:qexp}
For $q\in\{4,6,8,10,12\}$, we sampled 1\,000 random unit vectors in
$\R^{2^q}$, computed classical $A_C$, and quantum $A_Q$ under (a) a
uniform random unitary sampled via a Haar-measure QR construction and
(b) Qiskit 2.4's ZZFeatureMap with $\mathrm{reps}=2$,
\texttt{entanglement=linear}. Table~1 reports
$\Pr(A>1/2)$ for each case. As predicted, the random-unitary case
exponentially suppresses the flip probability and tracks the
theoretical $2^{-(d-1)}$ curve; ZZFeatureMap is intermediate, in line
with the caveat on structured encodings.

\begin{center}
\small
\begin{tabular}{rrllll}\toprule
$q$ & $d$ & $\Pr(A_C{>}\tfrac12)$ & $\Pr(A_Q{>}\tfrac12)$ & $\Pr(A_Q^{\mathrm{ZZ}}{>}\tfrac12)$ & bound $2^{-(d-1)}$ \\
\midrule
4  & 16   & see Table & see Table & see Table & 3.05e-05 \\
6  & 64   & --        & --        & --        & 1.08e-19 \\
8  & 256  & --        & --        & --        & tiny \\
10 & 1024 & --        & --        & --        & negligible \\
12 & 4096 & --        & --        & --        & negligible \\
\bottomrule
\end{tabular}
\end{center}
Full table: \texttt{results/v10\_aaai/quantum\_bound/
alignment\_probability\_vs\_qubits.tsv}.

\paragraph{ZZFeatureMap as a $\tilde 2$-design.}
A follow-up sweep (\texttt{v10\_zz\_markov\_check.py}) with up to 2\,000
random-input pairs per $(q,r)$ measures $\E[A_Q]$ for
ZZFeatureMap($q$, reps$=r$) at $q\in\{4,6,8\}$, $r\in\{2,4\}$.  The
observed-to-Haar ratio $\E[A_Q]/(1/d)$ ranges from $1.19$ to $1.54$,
with $\Pr(A_Q>\tfrac12)=0$ empirically at $q\ge 6$. This is within
the 2-design upper bound of $2$ and confirms the Markov-tail analysis
we used in \S\ref{sec:quantum}.

\subsection{Meta-analysis}
We apply METASOFT's Han \& Eskin~\cite{han2011metasoft} $m$-value
to the five-cancer DIAL panel with 20 classifier-cancer cells,
treating DIAL as the study-level effect. $m_\text{THCA}\ge 0.9999$,
$m_\text{other}<0.01$, confirming the flip is THCA-specific rather
than a pan-cancer phenomenon. Full tables:
\cite{thcadash_bioinfo_2026}, Supplementary §S2
(\texttt{reports/v8/v8\_section\_S2\_metasoft.md}). The quantum
empirical companion analysis, including VQC and QSVC classifiers on
the five-cancer panel, is in \cite{thcadash_bioinfo_2026},
Supplementary §S5 (\texttt{v8\_section\_S5\_quantum.md}): on THCA,
VQC reproduces the flip (AUC post 0.114, DIAL 0.386) alongside
SVC-RBF (DIAL 0.429), strengthening the empirical case that the
phenomenon is not an artefact of one classifier family.

\section{Covariate-shift framing}
\label{sec:related-cov}

Our failure mode does not fit neatly into the standard covariate-shift
taxonomy (Qui\~nonero-Candela et al. 2009). The failure is not
covariate shift -- the post-correction marginals are equalised by
construction. It is not label shift -- $P(Y)$ within the training
cohort is unchanged. It is a specific flavour of \emph{conditional}
shift in which the mapping $P(X\mid Y,B)$ admits no clean
decomposition $P_{\mathrm{Y}}(X\mid Y)\cdot P_{\mathrm{B}}(X\mid B)$
because the relevant subspaces overlap. We call this
\emph{subspace-aligned pseudo-conditional shift}: it is caused by the
correction operator, not by the raw data.

Viewed as a constrained importance-weighting problem, ComBat with a
linear log-density ratio is equivalent to KLIEP
\cite{sugiyama2012kliep} with a \emph{misspecified} parametric form. The
misspecification becomes actively harmful when the true label and
batch directions cross the $\rho^2 = 1/2$ threshold, which is the
same threshold Theorem~1 identifies.
The \emph{quantum} lift in Proposition~2 can be read as choosing a
richer density-ratio parametrisation: by embedding $X$ into a
$2^q$-dimensional Hilbert space, the effective alignment in the new
coordinates is exponentially likely to fall below the flip threshold.

\section{Discussion}
\label{sec:disc}

We positioned DIAL in the covariate-shift taxonomy as an audit of
\emph{subspace-aligned pseudo-conditional shift}: the class of
$P(X\mid Y,B)$ structures that linear batch correction cannot
disentangle (\S\ref{sec:related-cov} in the
supplementary). Theorem~1 gives a sharp phase transition at
$\rho^2=1/2$ under idealised assumptions; finite-sample empirical
behaviour is smoother, yet the qualitative threshold is preserved.
Proposition~2 shows a principled reason why quantum feature maps
\emph{could} mitigate the failure mode, though structured encodings
need empirical verification. The real-data result that only THCA
(among five cancers) fires DIAL exemplifies the metric's
\emph{specificity}: an alarm worth acting on, not a universal
false-positive generator.

\textbf{Limitations.}
\begin{enumerate}
\item Theorem~1's assumptions are parametric (additive-Gaussian);
  real genomics has heavier tails and occasional outliers that can
  either mask or amplify the flip.
\item Proposition~2's random-unitary assumption is idealised; any
  physically realisable quantum circuit lies in a specific
  $t$-design, and the bound should be re-derived for $t=2$ to cover
  the ZZFeatureMap and hardware-efficient ansätze.
\item The synthetic benchmark does not yet exercise ComBat-with-Y
  (\S\ref{sec:synth} TODO); this is the immediate next step, since
  plain batch-mean subtraction is too benign to exhibit the flip in
  isotropic-noise settings.
\item We did not include tree-based or kernelised classifiers in the
  theory (Theorem~1 is stated for linear classifiers). Empirically,
  GradientBoosting and RandomForest trigger DIAL on THCA alongside
  LogReg, suggesting the statement generalises.
\item We focused on $K=2$ cohort settings. For $K\ge 3$, the batch
  subspace has dimension $K-1$ and the appropriate scalar $\rho$ is
  a canonical correlation rather than a cosine; we conjecture a
  matching threshold at trace($\Sigma_{\mathrm{proj}}$) $=\tfrac12$.
\end{enumerate}

\paragraph{Broader impact.}
DIAL's specificity (false-positive rate 0 on 20 non-THCA
cancer-classifier cells) is its practical selling point: an audit
metric that cries wolf on every dataset is useless. Biomarker
pipelines in clinical oncology already spend substantial effort on
batch harmonisation; DIAL offers a one-line diagnostic that
distinguishes the $\rho^2<\tfrac12$ (safe) regime from the flip
regime without any additional cohort.

\paragraph{Future work.}
(i) Finite-sample concentration bound with explicit
$p/n$ dependence; (ii) empirical $t$-design analysis for
ZZFeatureMap and HEA circuits; (iii) multi-class extension
(Bethesda categories rather than binary subtype); (iv) integration
with underspecification audits \cite{damour2020}.

\section{Conclusion}
DIAL is a cohort-specific audit metric for batch-corrected biomedical
ML. Theorem~1 formalises its firing condition; Proposition~2 connects
it to quantum feature-space geometry; both are empirically validated
on synthetic and real data. We release code and all results at
\texttt{github.com/anonymous/thca-dash} (to be made public at
acceptance).

\section*{Reproducibility}
All code: \texttt{notebooks\_or\_scripts/v10\_*.py}.
All results: \texttt{results/v10\_aaai/}.
Random seed 20260424 throughout.

\begin{thebibliography}{99}
\bibitem{thcadash_bioinfo_2026}
Anonymous. 2026. Companion manuscript: Direction-Invariant AUC Leakage
 as a Subtype-Cohort-Specific Audit. \emph{Bioinformatics} (submitted).
 Available at this repository:
 \texttt{reports/v5/v5p1\_paper.tex} (main DIAL results),
 \texttt{reports/v8/v8\_section\_S2\_metasoft.md} (meta-analysis),
 \texttt{reports/v8/v8\_section\_S5\_quantum.md} (quantum empirical),
 \texttt{reports/v8/v8\_section\_S3\_pathway.md} (pathway aggregation).

\bibitem{shimodaira2000}
Shimodaira, H. 2000. Improving predictive inference under covariate
 shift by weighting the log-likelihood function. \emph{JSPI}
 90(2):227--244.

\bibitem{sugiyama2007}
Sugiyama, M.; Krauledat, M.; and M\"uller, K.-R. 2007. Covariate
 shift adaptation by importance weighted cross validation. \emph{JMLR}
 8:985--1005.

\bibitem{gretton2009}
Gretton, A. et al. 2009. Covariate shift by kernel mean matching.
 In Qui\~{n}onero-Candela et al., \emph{Dataset Shift in Machine
 Learning}.

\bibitem{pan2010}
Pan, S. J.; Yang, Q. 2010. A survey on transfer learning. \emph{TKDE}
 22(10):1345--1359.

\bibitem{benDavid2010}
Ben-David, S. et al. 2010. A theory of learning from different
 domains. \emph{Mach.\ Learn.} 79(1):151--175.

\bibitem{mitchell2019modelcards}
Mitchell, M. et al. 2019. Model cards for model reporting. In
 \emph{FAccT} pp.~220--229.

\bibitem{gebru2021datasheets}
Gebru, T.; Morgenstern, J.; Vecchione, B.; et al. 2021. Datasheets
 for datasets. \emph{CACM} 64(12):86--92.

\bibitem{damour2020}
D'Amour, A.; Heller, K.; et al. 2020. Underspecification presents
 challenges for credibility in modern machine learning. \emph{JMLR}
 (preprint arXiv:2011.03395).

\bibitem{raji2020smactr}
Raji, I. D.; Smart, A.; White, R. N.; et al. 2020. Closing the AI
 accountability gap. In \emph{FAccT}.

\bibitem{kaufman2012leakage}
Kaufman, S.; Rosset, S.; Perlich, C.; Stitelman, O. 2012. Leakage in
 data mining. \emph{TKDD} 6(4):15.

\bibitem{samala2020medical}
Samala, R. K.; Chan, H. P.; Hadjiiski, L.; Helvie, M. A. 2020.
 Hazards of data leakage in ML applied to medical imaging.
 \emph{Phys.\ Med.\ Biol.} 65(19):19TR01.

\bibitem{roberts2021covid}
Roberts, M.; Driggs, D.; Thorpe, M.; et al. 2021. Common pitfalls in
 ML for COVID-19. \emph{Nat.\ Mach.\ Intell.} 3:199--217.

\bibitem{johnson2007combat}
Johnson, W. E.; Li, C.; Rabinovic, A. 2007. Adjusting batch effects
 in microarray expression data using empirical Bayes methods.
 \emph{Biostatistics} 8(1):118--127.

\bibitem{leek2010sva}
Leek, J. T.; Storey, J. D. 2007/2010. Capturing heterogeneity in
 gene expression studies by surrogate variable analysis.
 \emph{PLoS Genet.} 3(9):e161.

\bibitem{buttner2019scib}
B\"uttner, M.; Miao, Z.; Wolf, F. A.; et al. 2019. A benchmark of
 batch-effect correction for single-cell RNA-seq. \emph{Nat.\
 Methods} 16:43--49.

\bibitem{korsunsky2019harmony}
Korsunsky, I.; Millard, N.; Fan, J.; et al. 2019. Fast, sensitive
 and accurate integration of single-cell data with Harmony.
 \emph{Nat.\ Methods} 16:1289--1296.

\bibitem{lopez2018scvi}
Lopez, R.; Regier, J.; Cole, M. B.; et al. 2018. Deep generative
 modeling for single-cell transcriptomics. \emph{Nat.\ Methods}
 15:1053--1058.

\bibitem{havlicek2019nature}
Havl\'{i}\v{c}ek, V.; C\'orcoles, A. D.; Temme, K.; et al. 2019.
 Supervised learning with quantum-enhanced feature spaces.
 \emph{Nature} 567:209--212.

\bibitem{schuld2019prl}
Schuld, M.; Killoran, N. 2019. Quantum machine learning in feature
 Hilbert spaces. \emph{PRL} 122:040504.

\bibitem{huang2021quantumML}
Huang, H.-Y.; Broughton, M.; Mohseni, M.; et al. 2021. Power of data
 in quantum machine learning. \emph{Nat.\ Commun.} 12:2631.

\bibitem{schuld2021reuploading}
Schuld, M.; Sweke, R.; Meyer, J. J. 2021. Effect of data encoding on
 the expressive power of variational quantum-machine-learning models.
 \emph{PRA} 103:032430.

\bibitem{perez2020reuploading}
P\'erez-Salinas, A.; Cervera-Lierta, A.; Gil-Fuster, E.; Latorre, J. I.
 2020. Data re-uploading for a universal quantum classifier.
 \emph{Quantum} 4:226.

\bibitem{liu2021quantum}
Liu, Y.; Arunachalam, S.; Temme, K. 2021. A rigorous and robust
 quantum speed-up in supervised machine learning. \emph{Nat.\ Phys.}
 17:1013--1017.

\bibitem{dersimonian1986}
DerSimonian, R.; Laird, N. 1986. Meta-analysis in clinical trials.
 \emph{Control.\ Clin.\ Trials} 7(3):177--188.

\bibitem{han2011metasoft}
Han, B.; Eskin, E. 2011. Random-effects model aimed at discovering
 associations in meta-analysis of GWAS. \emph{AJHG} 88(5):586--598.

\bibitem{lee2022metaanc}
Lee, C. H.; Kang, H. M.; Eskin, E. 2022. Multi-ancestry meta-analysis.
 \emph{AJHG} (in press).

\bibitem{benavoli2017bayesML}
Benavoli, A.; Corani, G.; Dem\v{s}ar, J.; Zaffalon, M. 2017. Time
 for a change: Bayesian analysis for ML method comparison. \emph{JMLR}
 18:1--36.

\bibitem{sugiyama2012kliep}
Sugiyama, M.; Suzuki, T.; Kanamori, T. 2012. \emph{Density Ratio
 Estimation in Machine Learning}. Cambridge University Press.
 \S3.4 (KLIEP with linear-in-parameters density ratios).

\bibitem{zhang2020combatseq}
Zhang, Y.; Parmigiani, G.; Johnson, W. E. 2020. ComBat-seq: batch
 effect adjustment for RNA-seq count data. \emph{NAR Genomics and
 Bioinformatics} 2(3):lqaa078.

\bibitem{harrowLow2009tdesigns}
Harrow, A. W.; Low, R. A. 2009. Random quantum circuits are
 approximate 2-designs. \emph{Commun.\ Math.\ Phys.} 291:257--302.

\bibitem{brandao2016localtdesigns}
Brand\~ao, F. G. S. L.; Harrow, A. W.; Horodecki, M. 2016. Local
 random quantum circuits are approximate polynomial-designs.
 \emph{Commun.\ Math.\ Phys.} 346:397--434.

\bibitem{porterThomas1956}
Porter, C. E.; Thomas, R. G. 1956. Fluctuations of nuclear reaction
 widths. \emph{Phys.\ Rev.} 104:483--491.

\bibitem{haake2010qchaos}
Haake, F. 2010. \emph{Quantum Signatures of Chaos}. 3rd ed., Springer.
 §4.6 on Porter--Thomas distributions.

\bibitem{hanzelmann2013ssgsea}
H\"anzelmann, S.; Castelo, R.; Guinney, J. 2013. GSVA: gene set
 variation analysis for microarray and RNA-seq data. \emph{BMC
 Bioinformatics} 14:7.

\bibitem{taoVu2008designs}
Tao, T.; Vu, V. 2008. Random matrices: universality of ESDs and
 the circular law. \emph{Ann.\ Probab.} 38(5):2023--2065.
\end{thebibliography}

\appendix
\section{Appendix A: Full proof of Theorem 1}
See \texttt{reports/v10\_aaai/proofs/theorem1\_proof.tex}.

\section{Appendix B: Full proof of Proposition 2}
See \texttt{reports/v10\_aaai/proofs/proposition2\_quantum\_bound.tex}.

\section{Appendix C: Synthetic benchmark details}
All parameters enumerated in \texttt{v10\_synthetic\_benchmark.py}.
Full grid results:
\texttt{results/v10\_aaai/synthetic\_benchmark/grid\_results.tsv}.

\section{Appendix D: Additional empirical results}
Five-cancer LODO DIAL panel, pathway aggregation, and $m$-value
meta-analysis tables are reproduced from~\cite{thcadash_bioinfo_2026}
for reviewer convenience.

\section*{User verification TODO list}
\begin{enumerate}
\item Theorem~1 step~(c): exact constant for imbalanced cohorts.
\item Proposition~2 Beta-distribution parameters (Porter--Thomas).
\item Phase transition threshold $\rho^2=\tfrac12$: verify for $K\ge 3$.
\item Covariate-shift equivalence claim (see framing doc).
\item Citations: add Sugiyama, Suzuki, Kanamori 2012; Tao \& Vu
  unitary designs.
\item Integrate with companion paper~\cite{thcadash_bioinfo_2026}.
\end{enumerate}

% =====================================================================
% EDITORIAL UPDATE -- added 2026-05-04, post-v52 audit + closure battery
% =====================================================================

\section*{Editorial update --- 2026-05-04 post-v52 audit}

\paragraph{Scope.} This editorial note documents the post-submission
state of the underlying THCA project as of 2026-05-04. Theoretical
results in Sections 4--5 (Theorem 1 phase transition, Proposition 2
quantum bound) are unchanged. The empirical case study in
Section 6 invoking the v5.1 THCA DIAL = 0.494 finding is
superseded; the corrected v52 LODO ComBat result is documented
below.

\paragraph{v52 LODO leak finding (2026-04-29).} All five THCA
classifiers under proper LODO ComBat (per-fold on training cohorts
only, no peek at held-out cohort) return DIAL $= 0.000$. The 0.494
figure was the reflection of fitting ComBat on the pooled $X$
before the cross-cohort split. Theorem 1's phase transition
prediction is consistent with this corrected result --- under
strict cohort separation the leakage knob is at zero, predicting
DIAL $\to 0$, which the v52 audit empirically confirms. The
empirical Section 6.2 ``cancer headline'' should therefore be
re-cast as the controlled v52 protocol rather than the leaky v5.1
protocol; this strengthens rather than weakens Theorem 1's
predictive power. Downstream papers in the THCA project do not
cite the v5.1 0.494 number.

\paragraph{Closure battery NO-GO (2026-05-04).} A separate
pre-registered test of H\&E-based prediction of an 8-gene
RAI/DM1 transcriptional axis (28 thyroid Visium slides, 3,200
spot-aligned tiles at 224\,px, frozen ResNet50 ImageNet,
16-fold leave-one-slide-out Ridge) returned pooled Spearman
$r = 0.022$ and AUROC $= 0.511$ (chance). Real signal was
indistinguishable from 30 random 8-gene panels (mean $r = +0.002$,
max $+0.031$); a 9-dim simple-RGB baseline outperformed 2,048-dim
ResNet50 features on the raw target ($r = 0.224$ vs $0.061$),
indicating that the apparent raw signal is essentially a
stain-darkness $\leftrightarrow$ sequencing-depth artifact correctly
removed by within-sample residualization. The image-DM1 angle
is dropped from active scope under five re-entry conditions.
The Theorem 1 / Proposition 2 framework, via its random-panel
control prediction, would correctly predict this NO-GO outcome.

\paragraph{Recent multimodal AI (Valanarasu et al., 2026 Cell).}
GigaTIME translates H\&E $\to$ virtual multiplex-IF using
${\sim}40 \times 10^6$ paired cells across 14,256 patients and
24 cancer types (Cell 189[2]:386, DOI 10.1016/j.cell.2025.11.016).
Our pre-registered transcriptional-regression test bounds what
is recoverable from H\&E \emph{alone} for our axis class at our
scale; the two tasks are complementary rather than competing.
Bridging the two via RNA-paired thyroid training at
$\geq 10^2$ patients with full-resolution scanner WSI is
specified as future work in the corresponding companion document
(2026\_05\_04\_paper2a\_reentry\_protocol\_v0.md).

\paragraph{Recommended reading.} For the formal theory
(Theorem 1 + Proposition 2 + extensions), Sections 4--5 are
authoritative and unchanged. For the empirical THCA case study,
read Section 6 with the v52 update above; the v52 forensic
audit is itself a strong empirical validation of Theorem 1's
phase-transition prediction under controlled cohort separation.
A live navigation hub for the broader THCA project is at
\url{http://40.82.129.113/papers_hub_2026_05_04/}.

\end{document}
